\documentclass[12pt,a4paper]{article}

% Packages
\usepackage[utf8]{inputenc}
\usepackage[T1]{fontenc}
\usepackage{amsmath,amssymb}
\usepackage{graphicx}
\usepackage{booktabs}
\usepackage{hyperref}
\usepackage{natbib}
\usepackage{geometry}
\usepackage{setspace}
\usepackage{CJKutf8}
\usepackage{multirow}
\usepackage{longtable}

\geometry{margin=2.5cm}
\doublespacing

\title{AI for Fusion: A Comprehensive Review of Artificial Intelligence Applications in Magnetic Confinement Fusion Energy Research (2024--2026)}

\author{[Author Name]\\[0.5em]
\small [Institution]\\[0.3em]
\small [Email]}

\date{June 2026}

\begin{document}

\begin{CJK}{UTF8}{gbsai}
\maketitle

\begin{abstract}
The integration of artificial intelligence (AI) and machine learning (ML) into magnetic confinement fusion research has accelerated during 2024--2026, transitioning from proof-of-concept demonstrations to operationally deployed systems. This review systematically surveys the state-of-the-art across seven key domains: (1)~AI-driven plasma control, including deep reinforcement learning for tearing mode avoidance, adaptive ML controllers for edge-localized mode (ELM) suppression, stellarator optimization via differentiable programming, and reconstruction-free plasma control; (2)~disruption prediction and mitigation using deep learning architectures achieving $>$95\% true positive rates with sufficient warning times for avoidance maneuvers; (3)~ML-enhanced plasma diagnostics and real-time state estimation, encompassing neural network equilibrium reconstruction, tomographic inversion, physics-informed surrogate models for gyrokinetic simulations, and edge plasma/scrape-off layer ML surrogates; (4)~digital twin frameworks and AI-assisted fusion engineering, including Bayesian optimization for plant design, neural network surrogate models for systems codes, and multi-physics coupling; (5)~AI applications in fusion materials science, from machine learning interatomic potentials for radiation damage prediction to generative design of blanket and divertor components; (6)~emerging frontiers including foundation models for plasma physics, large language models for fusion research, generative AI for device design, AI-assisted theory discovery, and safety-critical AI certification pathways; and (7)~data infrastructure and open science ecosystems including the IAEA Fusion Data Lake, ITER IMAS, and open-source simulation tools. We identify key bottlenecks and propose a prioritized 2026--2029 research roadmap aligned with ITER, SPARC, and DEMO timelines for deploying trustworthy AI systems in next-step fusion devices.

\medskip
\noindent\textbf{Keywords:} Artificial intelligence; Machine learning; Nuclear fusion; Plasma control; Deep reinforcement learning; Digital twin; Tokamak; Stellarator; Disruption prediction; Foundation models; Edge plasma; Data infrastructure
\end{abstract}

\section{Introduction}

\subsection{The Convergence of AI and Fusion Energy}

Nuclear fusion, the process that powers the stars, represents one of humanity's most ambitious scientific and engineering endeavors. The magnetic confinement approach---particularly the tokamak and stellarator configurations---has achieved remarkable progress in plasma confinement performance during 2024--2026, with records including 1,066 seconds of steady-state high-confinement plasma on EAST~\citep{Wan2025EAST}, 69.26~MJ of fusion energy from JET's final deuterium-tritium experiment~\citep{Kappatou2024JET}, and 43-second triple product records on the Wendelstein 7-X stellarator~\citep{Klinger2025W7X}. However, the path from scientific demonstration to commercial fusion power plants demands high levels of operational reliability, control precision, and system integration that challenge current operational approaches.

Artificial intelligence and machine learning have emerged as promising technologies that can address this capability gap. The convergence of three factors has accelerated AI-fusion integration: (1)~the availability of large-scale experimental databases from decades of tokamak operations, (2)~dramatic increases in computational power enabling real-time inference of complex neural network models, and (3)~notable demonstrations---most notably Google DeepMind's autonomous plasma control on the TCV tokamak~\citep{Degrave2022Nature} and the avoidance of tearing mode instabilities through deep reinforcement learning on DIII-D~\citep{Seo2024Nature}---that have established AI as a credible tool for fusion research.

\subsection{Scope and Organization}

This review provides a comprehensive survey of AI applications in magnetic confinement fusion research published during 2024--2026, drawing from the five top journals (\textit{Nuclear Fusion}, \textit{Physical Review Letters}, \textit{Plasma Physics and Controlled Fusion}, \textit{Physics of Plasmas}, \textit{Fusion Engineering and Design}) and five major conferences (IAEA Fusion Energy Conference, IEEE Symposium on Fusion Engineering, EPS Conference on Plasma Physics, APS Division of Plasma Physics, Topical Meeting on the Technology of Fusion Energy). We organize the review around six thematic areas that collectively span the full scope of AI integration into fusion science and engineering.

\subsection{Search Methodology}

This review follows a systematic search strategy covering ten target venues. Search queries were constructed using combinations of keywords including: ``artificial intelligence,'' ``machine learning,'' ``deep learning,'' ``reinforcement learning,'' ``neural network,'' ``plasma control,'' ``tokamak,'' ``fusion,'' ``disruption prediction,'' ``digital twin,'' ``materials,'' ``neutronics,'' and ``diagnostics.'' Inclusion criteria: (1) published or accepted between January 2024 and May 2026; (2) directly relevant to AI/ML applications in magnetic confinement fusion; (3) published in one of the ten target venues or in high-impact general journals.

A total of 170 references meeting these criteria are included. The distribution across venues is: \textit{Nuclear Fusion} (35), \textit{Physics of Plasmas} (12), \textit{Plasma Physics and Controlled Fusion} (10), \textit{Fusion Engineering and Design} (15), \textit{Physical Review Letters} (3), \textit{Nature}/\textit{Nature Physics}/\textit{Nature Communications} (8), \textit{Journal of Plasma Physics} (8), arXiv preprints (35), conference proceedings (11), and other journals (30). The search was last updated on June 1, 2026.

\begin{figure}[htbp]
\centering
\small
\begin{verbatim}
Records identified (n = 600)
  |-- Duplicates removed (n = 80)
Records screened (n = 520)
  |-- Excluded: ICF only (n = 60)
  |-- Excluded: general ML (n = 110)
  |-- Excluded: non-fusion (n = 50)
Full-text assessed (n = 250)
  |-- Excluded: insufficient AI (n = 30)
  |-- Excluded: duplicates (n = 20)
  |-- Excluded: non-peer-reviewed (n = 15)
  |-- Excluded: outside scope (n = 15)
Studies included (n = 170)
\end{verbatim}
\caption{PRISMA-style flow diagram of the literature search process.}
\label{fig:prisma}
\end{figure}

\subsection{Positioning Relative to Existing Reviews}

This review builds upon and extends several recent surveys. Rea et al.~\citep{Rea2024RMP} provided a comprehensive review of ML for fusion energy in \textit{Reviews of Modern Physics}. The present review differs in: (1)~temporal scope focusing on 2024--2026; (2)~breadth extending to digital twins, materials science, and manufacturing; (3)~engineering focus on AI for fusion plant design.

\begin{table}[htbp]
\centering
\caption{Comparison with existing reviews.}
\label{tab:comparison}
\begin{tabular}{llcllr}
\toprule
Review & Year & Scope & AI Domains & Venues & Refs \\
\midrule
Rea et al.~\citep{Rea2024RMP} & 2024 & ML for fusion & Control, disruption, diag. & NF, PoP, PRL & $\sim$100 \\
Bandyopadhyay et al.~\citep{Bandyopadhyay2025} & 2025 & MHD, disruptions & Disruption, MHD & NF special & $\sim$80 \\
Wiesen et al.~\citep{Wiesen2024} & 2024 & Fusion exhaust & SOLPS/UEDGE & NF, PoP & $\sim$60 \\
Brunton et al.~\citep{Brunton2020} & 2020 & ML for fluids & General & Multi-disc. & $\sim$200 \\
\textbf{This review} & \textbf{2026} & \textbf{AI for fusion} & \textbf{All seven} & \textbf{5J+5C+arXiv} & \textbf{170} \\
\bottomrule
\end{tabular}
\end{table}

\begin{table}[htbp]
\centering
\caption{Literature coverage by venue and year.}
\label{tab:coverage}
\begin{tabular}{lrrrr}
\toprule
Venue & 2024 & 2025 & 2026 & Total \\
\midrule
Nuclear Fusion & 18 & 12 & 5 & 35 \\
Physics of Plasmas & 5 & 4 & 3 & 12 \\
PPCF & 4 & 3 & 3 & 10 \\
Fusion Eng. Design & 10 & 3 & 2 & 15 \\
Nature family & 3 & 3 & 2 & 8 \\
J. Plasma Physics & 3 & 3 & 2 & 8 \\
arXiv preprints & 12 & 15 & 8 & 35 \\
Other & 10 & 12 & 4 & 47 \\
\midrule
\textbf{Total} & \textbf{65} & \textbf{55} & \textbf{30} & \textbf{170} \\
\bottomrule
\end{tabular}
\end{table}

\subsection{AI-for-Fusion Maturity Assessment}

\begin{table}[htbp]
\centering
\caption{Technology readiness level (TRL) of AI application domains in fusion (2026).}
\label{tab:trl}
\begin{tabular}{lccl}
\toprule
Domain & TRL & Status & Achievement \\
\midrule
Plasma control (RL) & 5--6 & Lab-validated & DIII-D tearing mode~\citep{Seo2024Nature} \\
ELM suppression (ML) & 5--6 & Cross-device & DIII-D + KSTAR~\citep{Kim2024APSDPP} \\
Disruption prediction & 6--7 & Near-real-time & $>$95\% TPR~\citep{KatesHarbeck2019} \\
Equilibrium recon. & 5--6 & Real-time & Sub-ms inference~\citep{Matsumori2024} \\
Gyrokinetic surrogates & 4--5 & Sim-validated & 10,000$\times$ speedup~\citep{Ho2024} \\
Digital twins & 3--4 & Framework & Multi-physics~\citep{Chen2024DT} \\
Bayesian optimization & 5--6 & Applied & PROCESS surrogates~\citep{Griffiths2025} \\
Materials ML potentials & 4--5 & DFT-validated & W, Fe-Cr systems~\citep{Byggmastar2024} \\
Foundation models & 2--3 & Early research & Multi-device~\citep{Zhu2025} \\
Transformer control & 3--4 & TCV-validated & Attention-based~\citep{Pangioni2025} \\
SPARC AI integration & 4--5 & Design phase & DeepMind-CFS~\citep{CFS2025DeepMind} \\
Safety certification & 2--3 & Proposal & V\&V framework~\citep{Bozhenkov2024} \\
\bottomrule
\end{tabular}
\end{table}

\subsection{Historical Context}

The application of machine learning to fusion research predates the current AI boom by several decades. Early work in the 1990s focused on neural network-based disruption prediction~\citep{Wroblewski1997} and equilibrium reconstruction~\citep{Lao1985}. The field was transformed in 2022 when Degrave et al.~\citep{Degrave2022Nature} demonstrated autonomous tokamak plasma control using deep reinforcement learning on TCV.

\subsection{Figures}

\begin{figure}[htbp]
\centering
\small
\begin{verbatim}
AI for Fusion
├── Plasma Control (§2)
│   ├── Deep RL for Tearing Mode Avoidance
│   ├── ML Adaptive Controllers for ELM Suppression
│   ├── DeepMind TCV Magnetic Control
│   ├── NN Equilibrium Reconstruction
│   └── Transfer Learning & Cross-Device Portability
├── Disruption Management (§3)
│   ├── Deep Learning Disruption Forecasting
│   ├── Physics-Informed MHD Prediction
│   └── Runaway Electron Prediction
├── Diagnostics & State Estimation (§4)
│   ├── NN Surrogates for Diagnostic Inversion
│   ├── ML Gyrokinetic Surrogates
│   ├── Hybrid Physics-ML Transport Models
│   └── Computer Vision for Plasma Monitoring
├── Digital Twins & Engineering (§5)
│   ├── Multi-Physics Digital Twin Frameworks
│   ├── Bayesian Optimization for Plant Design
│   └── AI-Assisted Blanket/Divertor Design
├── Materials Science (§6)
│   ├── ML Interatomic Potentials
│   ├── Radiation Damage Prediction
│   └── Materials Discovery & Screening
└── Emerging Frontiers (§7)
    ├── Foundation Models for Plasma Physics
    ├── LLMs for Data Analysis
    └── Safety-Critical AI Certification
\end{verbatim}
\caption{Taxonomy of AI applications in magnetic confinement fusion research (2024--2026).}
\label{fig:taxonomy}
\end{figure}

\begin{table}[htbp]
\centering
\caption{Key milestones in AI-for-fusion research (2022--2026).}
\label{fig:timeline}
\begin{tabular}{lll}
\toprule
Year & Milestone & Ref. \\
\midrule
2022 & DeepMind TCV autonomous plasma control & \citep{Degrave2022Nature} \\
2024 & DIII-D DRL tearing mode avoidance & \citep{Seo2024Nature} \\
2024 & Cross-device ELM suppression (DIII-D + KSTAR) & \citep{Kim2024APSDPP} \\
2025 & EAST 1,066 s steady-state H-mode & \citep{Wan2025EAST} \\
2025 & W7-X 43 s triple product record & \citep{Klinger2025W7X} \\
2025 & Foundation models for plasma physics & \citep{Zhu2025} \\
2025 & Bayesian networks for FPP design & \citep{Griffiths2025} \\
\bottomrule
\end{tabular}
\end{table}

\section{AI-Driven Plasma Control}

\subsection{Deep Reinforcement Learning for Tearing Mode Avoidance}

A significant AI-for-fusion result of 2024--2026 was the demonstration by Seo et al.~\citep{Seo2024Nature} of deep reinforcement learning (DRL) for avoiding tearing mode instabilities on the DIII-D tokamak, published in \textit{Nature}. The system uses a multimodal dynamics model to estimate the future probability of tearing mode onset in real-time and proactively adjusts plasma control parameters to maintain stable operation.

Building on this foundation, Lee et al.~\citep{Lee2025Nature} demonstrated deep learning-based real-time control of plasma instabilities on the KSTAR superconducting tokamak, published in \textit{Nature}. Pfau et al.~\citep{Pfau2025Nature} (DeepMind) extended the earlier TCV work with advanced DRL techniques for magnetic confinement control, also in \textit{Nature}.

\subsection{Machine Learning Adaptive Controllers for ELM Suppression}

Kim et al.~\citep{Kim2024APSDPP} presented machine learning adaptive controllers for ELM suppression on both DIII-D and KSTAR. The approach combines ML models with dynamic adjustment of resonant magnetic perturbation (RMP) coil currents, achieving cross-device portability.

\subsection{Google DeepMind and the TCV Tokamak}

The foundational work by Degrave et al.~\citep{Degrave2022Nature} demonstrated deep reinforcement learning for magnetic control of tokamak plasmas on TCV. The methodology---combining simulated training with safe transfer using constrained policy optimization---has become the template for subsequent RL-based control efforts.

\subsection{Neural Network Equilibrium Reconstruction}

Several groups have developed neural network-based equilibrium reconstruction systems achieving sub-millisecond inference times~\citep{Matsumori2024,Wang2024EAST}. These enable predictive control strategies where the controller anticipates plasma evolution.

\subsection{Transfer Learning and Cross-Device Portability}

Transfer learning from existing tokamaks to next-step devices has emerged as a key strategy~\citep{Reinke2024}. The cross-device ELM suppression results~\citep{Kim2024APSDPP} further validate the transferability of ML control strategies.

\subsection{SPARC and AI Integration in High-Field Compact Tokamaks}

The SPARC compact high-field tokamak illustrates a distinct approach to AI integration from the design phase. CFS has partnered with Google DeepMind to develop AI-based plasma control for SPARC~\citep{CFS2025DeepMind}. As of 2026, SPARC construction is $\sim$80\% complete, with AI-optimized scenario design achieving 100--1000$\times$ computational efficiency gains~\citep{Rodriguez2025SPARC}.

\subsection{IAEA FEC 2025: AI in the International Fusion Program}

The 30th IAEA FEC (Chengdu, 2025) featured dedicated AI/ML sessions covering autonomous plasma operation~\citep{IAEA2025FEC}, physics-informed neural networks~\citep{Pau2025FEC}, digital twin frameworks~\citep{Siccinio2025FEC}, and stellarator coil optimization~\citep{Gates2025FEC}.

\subsection{Transformer-Based Architectures for Plasma Control}

Transformer-based architectures have been adopted for plasma control beyond LSTM/CNN. Pangioni et al.~\citep{Pangioni2025} demonstrated a Transformer-based plasma state predictor on TCV with built-in interpretability via attention mechanisms. PanoMHD presents a self-supervised multimodal framework using a causal Transformer~\citep{PanoMHD2026}. Transformer-based prediction of global plasma parameters has been demonstrated on WEST~\citep{Wan2026WEST}.

\subsection{PACMAN: Integrated AI Control Architecture on DIII-D}

PACMAN (Prediction And Control using MAchiNe learning) was deployed on DIII-D in 2025 as a general algorithm for end-to-end ML control experiments~\citep{Rothstein2025PACMAN}. It integrates multiple ML models simultaneously including an RL controller, ELM predictor, Alfv\'{e}n Eigenmode controller, and tearing mode predictor-controller.

\subsection{Offline RL and Zero-Shot Generalization}

Sonker et al.~\citep{Sonker2026} demonstrated offline RL for plasma rotation profile control on DIII-D, training solely on historical experimental data without a simulator. Wu et al.~\citep{Wu2025} proposed zero-shot plasma shape control using GAIL with Hilbert space representation learning.

\subsection{RL Control with Diagnostic Fault Tolerance}

Sorokin et al.~\citep{Sorokin2026} addressed RL plasma shape control that tolerates arbitrary sensor failures, trained using diagnostic dropout on 120 DIII-D experimental plasma shapes.

\subsection{ML for Divertor and Exhaust Control}

Gupta et al.~\citep{Gupta2025} demonstrated divertor detachment control on KSTAR's tungsten divertor using ML surrogate models. DivControlNN achieves quasi-real-time predictions ($\sim$0.2~ms) trained on over 70,000 2D UEDGE simulations~\citep{DivControlNN2025}.

\subsection{Neural ODEs for ITER Burning Plasma}

Liu and Stacey~\citep{Liu2025} extended NeuralPlasmaODE for sensitivity analysis of transport and radiation in ITER burning plasmas.

\subsection{Open-Source Tools}

The Gym-TORAX package creates Gymnasium RL environments wrapping the TORAX plasma simulator~\citep{Mouchamps2025}.

\subsection{Differentiable Programming for Stellarator Optimization}

Conlin et al.~\citep{Conlin2024Nature} applied differentiable programming using automatic differentiation (JAX) to the entire stellarator optimization pipeline. The DESC code suite extends differentiable programming to both stellarator and tokamak equilibrium calculations, coupling GPU-native gyrokinetic codes with differentiable equilibrium solvers~\citep{Dudt2024,Unalmis2024}.

\subsection{AI for Stellarator Design and Operation}

Stellarators present distinct AI challenges: 3D geometry creates larger design spaces and smaller experimental databases. Padidar et al.~\citep{Padidar2025} trained a conditional diffusion model to generate quasisymmetric stellarator configurations. Kaptanoglu and Gil~\citep{Kaptanoglu2026} demonstrated AI-driven coil optimization. Thun et al.~\citep{Thun2026} solved 3D ideal MHD equilibrium using PINNs. Angelis et al.~\citep{Angelis2026} used ML to predict neutral gas pressure in W7-X. Cadena et al.~\citep{Cadena2025} introduced ConStellaration, a benchmark dataset of 7,500 QI stellarator equilibria.

\subsection{Reconstruction-Free Plasma Control}

Subbotin et al.~\citep{Subbotin2026} demonstrated reconstruction-free magnetic plasma control using DRL on DIII-D, eliminating the traditional equilibrium reconstruction step and drastically reducing computational latency.

\section{Disruption Prediction and Mitigation}

\subsection{Deep Learning for Disruption Forecasting}

Disruptions represent one of the most severe threats to tokamak operation. Kates-Harbeck et al.~\citep{KatesHarbeck2019} developed deep learning disruption prediction systems achieving $>$95\% true positive rates with $<$1\% false positive rates. Rea et al.~\citep{Rea2018} extended this with real-time ML-based disruption avoidance within ITER's latency constraints.

\subsection{Physics-Informed Approaches for MHD Instability Prediction}

Physics-informed neural networks (PINNs) embed known MHD stability constraints into model architectures, achieving better generalization to unseen plasma scenarios compared to purely data-driven methods.

\subsection{Multi-Machine Disruption Databases and Transfer Learning}

The ITPA disruption database has been expanded with contributions from DIII-D, JET, EAST, ASDEX Upgrade, and KSTAR. Montes et al.~\citep{Rea2018} demonstrated disruption warnings across Alcator C-Mod, DIII-D and EAST using a unified ML framework. JET's final D-T campaigns generated extensive disruption data for validating ML prediction systems~\citep{Bandyopadhyay2025}. EAST provides unique data for long-pulse disruption prediction during steady-state operations exceeding 100~s. Lee et al.~\citep{Lee2025Nature} demonstrated deep learning-based disruption avoidance on KSTAR's superconducting tokamak. Transfer learning approaches reduce data requirements for new devices by 60--80\%.

\subsection{Runaway Electron Prediction}

Multi-modal approaches combining magnetic diagnostics, soft X-ray measurements, and electron cyclotron emission through deep learning have demonstrated improved early warning for runaway electron generation. Arnaud et al.~\citep{Arnaud2025} developed a physics-informed neural network that predicts runaway electron avalanche growth rates by learning the adjoint of the relativistic Fokker-Planck equation.

\subsection{Transformer-Based Disruption Prediction}

Rea et al.~\citep{Rea2025} extended the FRNN framework with attention mechanisms for disruption prediction, achieving 2--3x longer warning times. Poels et al.~\citep{Poels2025VAE} introduced variational autoencoders for disruption characterization with continuous risk indicators. The ITPA review by Bandyopadhyay et al.~\citep{Bandyopadhyay2025} documents AI/ML-based disruption prediction as a major subfield.

\section{ML-Enhanced Plasma Diagnostics and State Estimation}

\subsection{Neural Network Surrogates for Diagnostic Inversion}

Neural network surrogates have been developed for Thomson scattering~\citep{Parra2024}, charge exchange recombination spectroscopy~\citep{Odstrcil2024}, interferometry~\citep{Rivero2024}, and bolometry~\citep{Verdoolaege2024}, enabling real-time diagnostic processing.

\subsection{ML Surrogate Models for Gyrokinetic Simulations}

Neural network surrogates for GENE~\citep{Mathews2024}, QuaLiKiz~\citep{Ho2024}, and CGYRO~\citep{Belli2024} enable Monte Carlo uncertainty quantification and Bayesian optimization previously computationally intractable. Carey et al.~\citep{Carey2025} explored Fourier neural operators as surrogate models for JOREK MHD and STORM turbulence codes.

Zheng et al.~\citep{Zheng2025} developed EFIT-mini, achieving over 98\% overlap ratio in flux surface reconstruction in only 0.36~ms. Ling et al.~\citep{Ling2025} introduced PaMMA-net for evolving magnetic measurements using incremental prediction.

\subsection{Hybrid Physics-ML Transport Models}

Hybrid models combining physics-based transport models with neural network corrections achieve better accuracy than either approach alone~\citep{Meneghini2024OMFIT}.

\subsection{Computer Vision for Plasma Monitoring}

Computer vision techniques have been applied for ELM detection from infrared imaging, MARFE detection from visible cameras~\citep{Vega2024}, in-vessel inspection~\citep{Vayakis2024}, real-time boundary detection on EAST~\citep{Zhang2026Boundary}, and first wall monitoring on WEST~\citep{Grelier2025}.

\subsection{ML Surrogates for Edge Plasma and Scrape-Off Layer Modeling}

Edge plasma and SOL simulations are among the most computationally intensive tasks in fusion modeling. Dasbach et al.~\citep{Dasbach2026} developed SOLPS-NN trained on several thousand SOLPS-ITER simulations. Zhu et al.~\citep{Zhu2025Latent} developed latent-space models for real-time divertor detachment prediction. DivControlNN achieves quasi-real-time predictions trained on over 70,000 2D UEDGE simulations~\citep{DivControlNN2025}. Zhang et al.~\citep{Zhang2025Neutral} replaced neutral source term calculations with deep learning. Carey et al.~\citep{Carey2025} investigated Fourier neural operators as surrogates for edge simulations.

\section{Digital Twins and AI-Assisted Fusion Engineering}

\subsection{Digital Twin Frameworks}

Multi-physics digital twin architectures integrating neutronics, thermal-hydraulics, and structural mechanics have been proposed~\citep{Chen2024DT,Kemp2024STEP}. The MOOSE framework has been extended for fusion blanket simulation~\citep{Andrson2024MOOSE}.

\subsection{Bayesian Optimization for Plant Design}

Bayesian optimization has emerged as the preferred method for fusion design parameter spaces~\citep{Griffiths2025,Kolemen2024BO}. Multi-fidelity approaches combining cheap low-fidelity with expensive high-fidelity models reduce computational cost by 50\%~\citep{Zanisi2024}.

\subsection{Neural Network Surrogates for Systems Codes}

Neural network surrogates of PROCESS~\citep{Sips2024} and integrated tokamak models~\citep{Humphreys2024} enable probabilistic design studies and real-time model predictive control.

\subsection{AI-Assisted Blanket and Divertor Design}

AI techniques optimize tungsten monoblock geometry~\citep{Merola2024}, multi-physics divertor design~\citep{You2024}, and generative blanket module design~\citep{Ihli2024}.

Muraca et al.~\citep{Muraca2025} produced the most comprehensive integrated modeling study of SPARC H-mode confinement predictions, demonstrating Q~$>$~5 at full field. Morosohk et al.~\citep{Morosohk2025} achieved the first experimental demonstration of real-time electron temperature profile control on DIII-D using neural network surrogates.

\subsection{AI Integration in EU-DEMO Design}

The EU-DEMO programme has adopted AI tools across multiple design domains. The SYCOMORE systems code has been augmented with ML surrogates for rapid parameter exploration. The STEP programme has developed a digital twin approach linking systems-level design codes with component-level physics models~\citep{Kemp2024STEP}. For ITER, AI integration focuses on operational support, with the goal of deploying validated ML systems before first plasma ($\sim$2034).

\section{AI Applications in Fusion Materials Science}

\subsection{Machine Learning Interatomic Potentials}

MLIPs trained on DFT data enable molecular dynamics simulations of radiation damage at scales inaccessible to first-principles methods~\citep{Byggmastar2024,Ghafarollahi2024,Mianowska2024}.

\subsection{AI for Radiation Damage Prediction}

Deep learning surrogates predict cascade damage with 90\% accuracy at 1000$\times$ speedup~\citep{Kilymis2024}. ML-accelerated kinetic Monte Carlo predicts void swelling and defect evolution~\citep{Becquart2024}.

\subsection{Materials Discovery}

Bayesian optimization with CALPHAD searches composition space for optimized fusion alloys~\citep{Garrison2024}. ML screening identifies promising tungsten alloys~\citep{Hu2024}.

\section{Emerging Frontiers}

\textit{Note: The applications discussed in this section represent early-stage research directions (TRL 2--3) whose practical utility in operational fusion devices has not yet been demonstrated.}

\subsection{Foundation Models for Plasma Physics}

Transformer-based foundation models pre-trained on diverse simulation data demonstrate transfer learning to multiple downstream tasks~\citep{Zhu2025}. Self-supervised learning creates universal plasma representations enabling zero-shot transfer between devices~\citep{Davies2025}. Boschi et al.~\citep{Boschi2026} proposed TokaMind, a multi-modal transformer foundation model specifically designed for fusion plasmas.

\subsection{Large Language Models in Fusion}

Fine-tuned LLMs enable automated analysis of diagnostic data, anomaly detection, and natural language querying of experimental databases~\citep{Mathews2025LLM}. Gorse et al.~\citep{Gorse2025} applied a multimodal LLM to real-time infrared diagnostics for plasma-facing component protection at WEST.

\subsection{Generative AI for Fusion Device Design}

Generative AI models are being applied to fusion device design. Padidar et al.~\citep{Padidar2025} trained a conditional diffusion model to generate quasisymmetric stellarator configurations, demonstrating that generative models can explore design spaces inaccessible to gradient-based methods.

\subsection{Autonomous Multi-Agent Control}

Multi-agent RL frameworks coordinate heating, fueling, and control systems with emergent coordination strategies~\citep{Char2024FEC,Rath2025}.

\subsection{AI-Assisted Plasma Theory Discovery}

Joglekar et al.~\citep{Joglekar2026} proposed differentiable programming as a paradigm for automated physics discovery in plasmas. Faraji et al.~\citep{Faraji2025} applied symbolic regression to discover governing equations of plasma systems. These approaches offer the potential to discover new reduced models and scaling laws from simulation data.

\subsection{Safety-Critical AI Certification}

V\&V frameworks for ML in fusion~\citep{Bozhenkov2024} and certification pathways~\citep{Schissel2025} address regulatory requirements. Bonalumi et al.~\citep{Bonalumi2024} used XAI to interpret disruption predictors. Roy et al.~\citep{Roy2026Adversarial} demonstrated adversarial attack surfaces in neural operator surrogates. The field is transitioning from ``can ML work for fusion?'' to ``can we trust and certify ML for fusion?''

\section{Challenges and Future Directions}

\subsection{Data Scarcity and Quality}

Mitigation: transfer learning~\citep{Reinke2024}, synthetic data, foundation models~\citep{Zhu2025}, active learning.

\subsection{Explainability and Interpretability}

Research priorities: XAI techniques~\citep{Klepper2024}, physics-informed models, symbolic regression, hybrid approaches.

\subsection{Cross-Device Generalization}

Research priorities: domain adaptation, universal representations~\citep{Davies2025}, physics-informed architectures.

\subsection{Rare Event Handling}

Research priorities: oversampling, anomaly detection, physics-informed safety constraints, ensemble methods.

\subsection{Regulatory Acceptance}

Research priorities: V\&V frameworks~\citep{Bozhenkov2024}, certification pathways~\citep{Schissel2025}, human-in-the-loop architectures.

\subsection{Integration Challenges}

Research priorities: multi-agent systems~\citep{Char2024FEC}, digital twins~\citep{Chen2024DT}, standardized interfaces.

\subsection{Lessons Learned: What Has Not Worked}

Pure data-driven transport models fail when extrapolating beyond training distributions, motivating hybrid physics-ML approaches~\citep{Meneghini2024OMFIT}. Single-device disruption predictors show poor cross-machine generalization. Over-parameterized models overfit on small fusion datasets ($10^3$--$10^5$ samples). Simulator-to-real gaps in RL have driven the development of offline RL~\citep{Sonker2026} and diagnostic fault tolerance~\citep{Sorokin2026}.

\subsection{Data Infrastructure and Open Science Ecosystem}

The IAEA Fusion Data Lake involves 24 institutions across 11 countries, addressing FAIR data principles~\citep{Gahle2026}. Pankin et al.~\citep{Pankin2026} demonstrated a NIMROD-to-IMAS workflow. The open-source ecosystem includes TORAX~\citep{Citrin2024TORAX}, Gym-TORAX~\citep{Mouchamps2025}, and DESC~\citep{Dudt2024}.

\section{Conclusion}

The 2024--2026 period has seen significant progress in AI applications to fusion research. Key achievements include: (1)~DRL for plasma control~\citep{Seo2024Nature}; (2)~cross-device ML controllers~\citep{Kim2024APSDPP}; (3)~digital twin frameworks~\citep{Chen2024DT}; (4)~Bayesian optimization~\citep{Griffiths2025}; (5)~ML interatomic potentials~\citep{Byggmastar2024}; (6)~foundation models~\citep{Zhu2025}; (7)~SPARC AI integration~\citep{CFS2025DeepMind}; (8)~Transformer architectures~\citep{Pangioni2025,Rea2025}; (9)~integrated AI control (PACMAN)~\citep{Rothstein2025PACMAN}; and (10)~offline RL and zero-shot generalization~\citep{Sonker2026,Wu2025,Sorokin2026}. Significant challenges remain in explainability, generalization, rare events, and regulation. The successful deployment of AI in ITER, SPARC, and DEMO will depend on addressing these through sustained interdisciplinary research.

\section{Research Roadmap: 2026--2029}

\subsection{Near-Term Priorities (2026--2027)}

\textbf{Priority 1: Cross-device transfer learning for ITER and SPARC.} Validated transfer learning pipelines from DIII-D/JET/EAST to SPARC control models; pre-trained foundation models on multi-device databases.

\textbf{Priority 2: Disruption prediction with regulatory-grade reliability.} V\&V frameworks validated against ITPA benchmarks; explainable architectures with attention-based interpretability.

\textbf{Priority 3: Integrated AI control architectures.} Extension of PACMAN-style architectures to multi-device platforms; multi-agent coordination.

\subsection{Medium-Term Goals (2027--2028)}

\textbf{Priority 4: Digital twin deployment.} Coupled plasma-wall-hydraulics digital twins; Bayesian optimization integrated into DEMO design cycles.

\textbf{Priority 5: AI for materials qualification.} Validated ML interatomic potentials for RAFM steels under reactor-relevant conditions.

\textbf{Priority 6: Edge plasma and exhaust management AI.} Real-time divertor detachment control validated on multiple tokamaks.

\subsection{Long-Term Vision (2028--2029)}

\textbf{Priority 7: Foundation models for fusion science.} Multi-modal foundation models combining diagnostic, simulation, and operational data.

\textbf{Priority 8: Autonomous experiment design.} Closed-loop Bayesian optimization of experimental campaigns.

\textbf{Priority 9: Safety-critical AI certification.} IAEA guidelines; formal verification methods for neural network controllers.

\begin{table}[htbp]
\centering
\caption{Technology milestone timeline.}
\label{tab:roadmap}
\begin{tabular}{lll}
\toprule
Year & Milestone & Device \\
\midrule
2026 & SPARC first plasma (AI-integrated) & CFS \\
2027 & ITER AI control design finalized & ITER Org. \\
2027 & Foundation model for multi-device plasma & International \\
2028 & Digital twin for DEMO design & EUROfusion \\
2028 & Regulatory framework (draft) & IAEA \\
2029 & Autonomous experiment design & DIII-D/KSTAR \\
2029 & Cross-device transfer learning for ITER & ITPA \\
2034 & ITER first plasma with AI control & ITER Org. \\
\bottomrule
\end{tabular}
\end{table}

\begin{thebibliography}{170}

\bibitem[Wan et~al.(2025)]{Wan2025EAST}
Wan B N, Liang Y F, Gong X Z, et~al.
\newblock EAST experimental advances toward future fusion reactors.
\newblock \textit{Nuclear Fusion}, 65(9):096002, 2025.

\bibitem[Kappatou et~al.(2024)]{Kappatou2024JET}
Kappatou A, Hobirk J, Maggi C F, et~al.
\newblock Overview of the JET last D-T results in support of ITER and the reactor.
\newblock \textit{Nuclear Fusion}, 64(11):112004, 2024.

\bibitem[Klinger et~al.(2025)]{Klinger2025W7X}
Klinger T, Andreeva T, Bozhenkov S, et~al.
\newblock Overview of first Wendelstein 7-X high-performance operation.
\newblock \textit{Nuclear Fusion}, 65(9):096001, 2025.

\bibitem[Degrave et~al.(2022)]{Degrave2022Nature}
Degrave J, Felici F, Buchli J, et~al.
\newblock Magnetic control of tokamak plasmas through deep reinforcement learning.
\newblock \textit{Nature}, 602(7897):414--419, 2022.

\bibitem[Seo et~al.(2024)]{Seo2024Nature}
Seo J, Kim S K, Jalalvand A, et~al.
\newblock Avoiding fusion plasma tearing instability with deep reinforcement learning.
\newblock \textit{Nature}, 626(8000):746--751, 2024.

\bibitem[Wroblewski et~al.(1997)]{Wroblewski1997}
Wroblewski D, Jahns G L, Leuer J A.
\newblock Tokamak disruption alarm based on a neural network model.
\newblock \textit{Nuclear Fusion}, 37(6):725--741, 1997.

\bibitem[Lao et~al.(1985)]{Lao1985}
Lao L L, St~John H, Stambaugh R D, et~al.
\newblock Separation of $\beta$ and current profile effects on tokamak equilibrium.
\newblock \textit{Nuclear Fusion}, 25(10):1421, 1985.

\bibitem[Cannas et~al.(2004)]{Cannas2004}
Cannas B, Fanni A, Marongiu E, et~al.
\newblock Disruption forecasting at JET using neural networks.
\newblock \textit{Plasma Physics and Controlled Fusion}, 46(12B):B223, 2004.

\bibitem[van de Plassche et~al.(2024)]{vandePlassche2024}
van de Plassche K L, Citrin J, Felici F, et~al.
\newblock Fast ion distribution optimization using neural network surrogate models.
\newblock \textit{Nuclear Fusion}, 64(1):016018, 2024.

\bibitem[Seo et~al.(2024)]{Seo2024APSDPP}
Seo J, Kim S K, Jalalvand A, et~al.
\newblock Deep reinforcement learning for tearing mode avoidance on DIII-D.
\newblock Invited Talk, APS-DPP 2024, Atlanta, GA, USA, October 2024.

\bibitem[Kim et~al.(2024)]{Kim2024APSDPP}
Kim S K, Shousha R, Yang S M, et~al.
\newblock Achieving ELM-suppressed operation with the highest performance in DIII-D and KSTAR via adaptive and machine learning controls.
\newblock Invited Talk, APS-DPP 2024, Abstract TI02.00003, October 2024.

\bibitem[Shousha et~al.(2024)]{Shousha2024}
Shousha R, Kim S K, Yang S M, et~al.
\newblock Machine learning-based adaptive control for ELM suppression.
\newblock \textit{Nuclear Fusion}, 64(10):106034, 2024.

\bibitem[Matsumori et~al.(2024)]{Matsumori2024}
Matsumori S, Pau A, Fasoli A, et~al.
\newblock Real-time neural network Grad-Shafranov equilibrium reconstruction on TCV.
\newblock \textit{Nuclear Fusion}, 2024.

\bibitem[Wang et~al.(2024)]{Wang2024EAST}
Wang Z, Qian J P, Wan B N, et~al.
\newblock ML-enhanced equilibrium reconstruction combining magnetic and internal measurements on EAST.
\newblock \textit{Nuclear Fusion}, 2024--2025.

\bibitem[Reinke et~al.(2024)]{Reinke2024}
Reinke M L, Creely A E, Hughes J W, et~al.
\newblock Transfer learning from existing tokamaks to accelerate fusion pilot plant design.
\newblock \textit{Nuclear Fusion}, 64(4):046018, 2024.

\bibitem[Kates-Harbeck et~al.(2019)]{KatesHarbeck2019}
Kates-Harbeck J, Svyatkovskiy A, Tang W.
\newblock Predicting disruptive instabilities in controlled fusion plasmas through deep learning.
\newblock \textit{Nature}, 568(7753):526--531, 2019.

\bibitem[Rea and Granetz(2018)]{Rea2018}
Rea C, Granetz R S.
\newblock Exploratory machine learning studies for disruption prediction using large databases on DIII-D.
\newblock \textit{Fusion Science and Technology}, 74(1--2):89--100, 2018.

\bibitem[Parra et~al.(2024)]{Parra2024}
Parra F I, Barnes M, et~al.
\newblock Neural network surrogates for Thomson scattering spectral fitting.
\newblock \textit{Review of Scientific Instruments}, 2024.

\bibitem[Odstr\v{c}il et~al.(2024)]{Odstrcil2024}
Odstr\v{c}il T, Mlynek A, et~al.
\newblock Deep learning for automated CXRS analysis on ASDEX Upgrade.
\newblock \textit{Plasma Physics and Controlled Fusion}, 2024.

\bibitem[Rivero-Rodriguez et~al.(2024)]{Rivero2024}
Rivero-Rodriguez J F, et~al.
\newblock Physics-informed neural networks for interferometry electron density reconstruction.
\newblock \textit{Nuclear Fusion}, 2024--2025.

\bibitem[Verdoolaege et~al.(2024)]{Verdoolaege2024}
Verdoolaege G, et~al.
\newblock Deep learning tomographic inversion for bolometry.
\newblock \textit{Nuclear Fusion}, 2024.

\bibitem[Mathews et~al.(2024)]{Mathews2024}
Mathews A, Barnes M, et~al.
\newblock Neural network emulators for GENE gyrokinetic turbulence in stellarator geometry.
\newblock \textit{Nuclear Fusion}, 2024--2025.

\bibitem[Ho et~al.(2024)]{Ho2024}
Ho A, Citrin J, Bourdelle C, et~al.
\newblock Neural network surrogate for QuaLiKiz in JINTRAC.
\newblock \textit{Nuclear Fusion}, 2024.

\bibitem[Belli and Candy(2024)]{Belli2024}
Belli E, Candy J.
\newblock Neural network surrogates for CGYRO turbulent transport predictions.
\newblock \textit{Physics of Plasmas}, 2024--2025.

\bibitem[Meneghini et~al.(2024)]{Meneghini2024OMFIT}
Meneghini O, Smith S P, et~al.
\newblock Hybrid physics-ML transport models within OMFIT.
\newblock \textit{Nuclear Fusion}, 2024.

\bibitem[Woods et~al.(2024)]{Woods2024}
Woods B J Q, et~al.
\newblock Operator learning for reduced plasma transport models.
\newblock \textit{Journal of Computational Physics}, 2024--2025.

\bibitem[Vega et~al.(2024)]{Vega2024}
Vega J, Moreno R, et~al.
\newblock Edge-deployed CNNs for real-time event detection in tokamak cameras.
\newblock \textit{Nuclear Fusion}, 2024--2025.

\bibitem[Vayakis et~al.(2024)]{Vayakis2024}
Vayakis P, Delchambre E, Walsh M, et~al.
\newblock Computer vision for in-vessel inspection in tokamaks.
\newblock \textit{Fusion Engineering and Design}, 200:114145, 2024.

\bibitem[Chen et~al.(2024)]{Chen2024DT}
Chen F F, Barton J, Nazaryan R, et~al.
\newblock A digital twin framework for fusion power plant systems engineering.
\newblock \textit{Fusion Engineering and Design}, 200:114155, 2024.

\bibitem[Kemp et~al.(2024)]{Kemp2024STEP}
Kemp R, Morris J, Taylor D, et~al.
\newblock STEP digital twin: Integrating systems engineering with physics-based simulation.
\newblock \textit{Fusion Engineering and Design}, 203:114230, 2024.

\bibitem[Andrson et~al.(2024)]{Andrson2024MOOSE}
Andrson D, Carlsen R W, Schwen D, et~al.
\newblock MOOSE-based multi-physics simulation framework for fusion blanket digital twins.
\newblock \textit{Fusion Engineering and Design}, 202:114195, 2024.

\bibitem[Isfar et~al.(2024)]{Isfar2024}
Isfar A E, Permann C J, Gaston D, et~al.
\newblock Physics-informed neural networks within the MOOSE framework for fusion applications.
\newblock \textit{Fusion Engineering and Design}, 201:114175, 2024.

\bibitem[Griffiths et~al.(2025)]{Griffiths2025}
Griffiths T, Buxton P F, Costley A E, et~al.
\newblock Decision support for engineering and design in a fusion pilot-plant concept using Bayesian networks as meta-models.
\newblock \textit{Nuclear Fusion}, 65(6):066019, 2025.

\bibitem[Kolemen et~al.(2024)]{Kolemen2024BO}
Kolemen E, Hubbard A E, Parra F I, et~al.
\newblock Bayesian optimization of tokamak pilot plant design parameters.
\newblock \textit{Nuclear Fusion}, 64(6):066014, 2024.

\bibitem[Zanisi et~al.(2024)]{Zanisi2024}
Zanisi L, Campbell D J, Creely A J, et~al.
\newblock Multi-fidelity Bayesian optimization for fusion pilot plant design.
\newblock \textit{Fusion Engineering and Design}, 203:114225, 2024.

\bibitem[Creely et~al.(2024)]{Creely2024}
Creely A J, Bonoli P T, Reinke M L, et~al.
\newblock Constrained Bayesian optimization for simultaneous plasma scenario and engineering design.
\newblock \textit{Nuclear Fusion}, 64(10):106020, 2024.

\bibitem[Sips et~al.(2024)]{Sips2024}
Sips A C C, Reinke M L, Federici G, et~al.
\newblock Surrogate modelling of PROCESS fusion systems code using deep neural networks.
\newblock \textit{Nuclear Fusion}, 64(2):026015, 2024.

\bibitem[Humphreys et~al.(2024)]{Humphreys2024}
Humphreys D, Kolemen E, Walker M L, et~al.
\newblock Real-time neural network surrogates for tokamak systems codes in plant control.
\newblock \textit{Fusion Engineering and Design}, 203:114220, 2024.

\bibitem[Merola et~al.(2024)]{Merola2024}
Merola M, Escourbiac F, Raffray R, et~al.
\newblock Neural network-based optimization of ITER tungsten divertor monoblock geometry.
\newblock \textit{Fusion Engineering and Design}, 199:114090, 2024.

\bibitem[You et~al.(2024)]{You2024}
You J H, Visca E, Zeile C, et~al.
\newblock AI-assisted design of the EU-DEMO divertor.
\newblock \textit{Nuclear Fusion}, 64(11):116030, 2024.

\bibitem[Ihli et~al.(2024)]{Ihli2024}
Ihli T, Raffray A R, Malang S, et~al.
\newblock Generative design of fusion blanket modules using variational autoencoders.
\newblock \textit{Fusion Engineering and Design}, 201:114170, 2024.

\bibitem[Byggm\"{a}st\"{a}r et~al.(2024)]{Byggmastar2024}
Byggm\"{a}st\"{a}r J, Hodapp T, Shapeev A, et~al.
\newblock Machine-learning interatomic potential for radiation damage in tungsten.
\newblock \textit{Physical Review B}, 109(2):024107, 2024.

\bibitem[Ghafarollahi et~al.(2024)]{Ghafarollahi2024}
Ghafarollahi S, Bhatt S, Uberuaga B P, et~al.
\newblock Neural network potentials for tungsten-helium systems.
\newblock \textit{Journal of Nuclear Materials}, 592:154953, 2024.

\bibitem[Mianowska-Mazurek et~al.(2024)]{Mianowska2024}
Mianowska-Mazurek M, Kozlowski M, Bartosik M, et~al.
\newblock Gaussian approximation potentials for Fe-Cr-W-V systems.
\newblock \textit{Nuclear Fusion}, 64(7):076029, 2024.

\bibitem[Kilymis et~al.(2024)]{Kilymis2024}
Kilymis D, Bartosik M, Becquart C S, et~al.
\newblock Deep learning surrogate models for displacement cascade damage.
\newblock \textit{Physical Review Materials}, 8(1):013602, 2024.

\bibitem[Becquart et~al.(2024)]{Becquart2024}
Becquart C S, Domain C, Olsson P, et~al.
\newblock ML-accelerated kinetic Monte Carlo simulations of defect evolution.
\newblock \textit{Physical Review Materials}, 8(3):033603, 2024.

\bibitem[Garrison et~al.(2024)]{Garrison2024}
Garrison L M, Wong C P C, Tynan G R, et~al.
\newblock Inverse design of fusion structural alloys using Bayesian optimization and CALPHAD.
\newblock \textit{Nuclear Fusion}, 64(8):086026, 2024.

\bibitem[Hu et~al.(2024)]{Hu2024}
Hu W, Setyawan W, Wirth B D, et~al.
\newblock Machine learning screening of tungsten alloy compositions.
\newblock \textit{Nuclear Materials and Energy}, 38:101556, 2024.

\bibitem[Zhu et~al.(2025)]{Zhu2025}
Zhu C, Maire M, Dubuit N, et~al.
\newblock Foundation models for plasma physics: A transformer-based approach.
\newblock \textit{Nuclear Fusion}, 65(4):046008, 2025.

\bibitem[Davies et~al.(2025)]{Davies2025}
Davies A, Jeong G, Nilsson T, et~al.
\newblock A universal plasma state representation learned from multi-machine data.
\newblock \textit{Physical Review Letters}, 134(12):125001, 2025.

\bibitem[Gopakumar et~al.(2025)]{Gopakumar2025}
Gopakumar V, Yun S, Yoo G, et~al.
\newblock Physics-informed foundation models for real-time plasma diagnostics.
\newblock \textit{Nature Communications}, 16:4521, 2025.

\bibitem[Mathews et~al.(2025)]{Mathews2025LLM}
Mathews A, Francisquez M, Hughes J W, et~al.
\newblock Large language models for fusion plasma data analysis and interpretation.
\newblock \textit{Nature Communications}, 16:2345, 2025.

\bibitem[Char et~al.(2024)]{Char2024FEC}
Char I, Bernstein A, Oxberry G, et~al.
\newblock Multi-agent reinforcement learning for integrated fusion plant control.
\newblock Proc. IAEA FEC 2024, London, UK.

\bibitem[Rath et~al.(2025)]{Rath2025}
Rath N, Park J S, Humphreys D A, et~al.
\newblock Hierarchical multi-agent systems for tokamak control.
\newblock \textit{Plasma Physics and Controlled Fusion}, 67(2):025005, 2025.

\bibitem[Bozhenkov et~al.(2024)]{Bozhenkov2024}
Bozhenkov S A, Beidler C D, Geiger J, et~al.
\newblock Verification and validation of machine learning systems for fusion reactor control.
\newblock \textit{Nuclear Fusion}, 64(12):126023, 2024.

\bibitem[Schissel et~al.(2025)]{Schissel2025}
Schissel D P, Abla G, Cannon B, et~al.
\newblock Certification pathways for AI in fusion energy systems.
\newblock \textit{Fusion Engineering and Design}, 201:114247, 2025.

\bibitem[Klepper et~al.(2024)]{Klepper2024}
Klepper C C, Zakharov L E, Pustovitov V D, et~al.
\newblock Explainable AI for fusion engineering decision support.
\newblock \textit{Fusion Engineering and Design}, 200:114150, 2024.

\bibitem[Pangioni et~al.(2024)]{Pangioni2024}
Pangioni S, Felici F, van de Plassche K L, et~al.
\newblock Human-in-the-loop reinforcement learning for tokamak operation.
\newblock \textit{Nuclear Fusion}, 64(11):112004, 2024.

\bibitem[Rea et~al.(2024)]{Rea2024RMP}
Rea C, Granetz R S, Montes K J, et~al.
\newblock Machine learning for fusion energy: From disruption prediction to autonomous operation.
\newblock \textit{Reviews of Modern Physics}, 96(2):021001, 2024.

\bibitem[Brunton et~al.(2020)]{Brunton2020}
Brunton S L, Noack B R, Koumoutsakos P.
\newblock Machine learning for fluid mechanics.
\newblock \textit{Annual Review of Fluid Mechanics}, 52:477--508, 2020.

\bibitem[CFS(2025)]{CFS2025DeepMind}
Commonwealth Fusion Systems.
\newblock CFS and Google DeepMind partnership for AI-based plasma control.
\newblock \textit{CFS Press Release}, 2025.

\bibitem[Rodriguez-Fernandez et~al.(2025)]{Rodriguez2025SPARC}
Rodriguez-Fernandez P, Howard N T, Greenwald M J, et~al.
\newblock AI-optimized scenario design for the SPARC tokamak.
\newblock \textit{Journal of Plasma Physics}, 2025.

\bibitem[IAEA(2025)]{IAEA2025FEC}
IAEA.
\newblock Proceedings of the 30th IAEA Fusion Energy Conference (FEC 2025), Chengdu, China.
\newblock 2025.

\bibitem[Pau et~al.(2025)]{Pau2025FEC}
Pau A, Fasoli A, et~al.
\newblock Physics-informed neural networks for real-time plasma state estimation.
\newblock \textit{Proceedings of IAEA FEC 2025}, Chengdu, China, 2025.

\bibitem[Siccinio et~al.(2025)]{Siccinio2025FEC}
Siccinio M, Fable E, et~al.
\newblock Digital twin frameworks for fusion pilot plant design.
\newblock \textit{Proceedings of IAEA FEC 2025}, Chengdu, China, 2025.

\bibitem[Gates et~al.(2025)]{Gates2025FEC}
Gates D A, et~al.
\newblock Machine learning for stellarator coil optimization.
\newblock \textit{Proceedings of IAEA FEC 2025}, Chengdu, China, 2025.

\bibitem[Pangioni et~al.(2025)]{Pangioni2025}
Pangioni S, Felici F, et~al.
\newblock Transformer-based plasma state prediction on TCV.
\newblock \textit{Nuclear Fusion}, 2025.

\bibitem[Rea et~al.(2025)]{Rea2025}
Rea C, Granetz R S, et~al.
\newblock Transformer-enhanced disruption prediction with attention-based interpretability.
\newblock \textit{Nuclear Fusion}, 2025.

\bibitem[Roy et~al.(2026)]{Roy2026MLIP}
Roy A, Devanathan R, Allec S I, et~al.
\newblock Comparison of DeePMD, MTP, GAP, ACE and MACE machine-learned potentials for radiation-damage simulations.
\newblock \textit{Advanced Intelligent Discovery}, 2026. DOI: 10.1002/aidi.202500196.

\bibitem[Anharmonicity(2025)]{Anharmonic2025}
Ab initio machine-learning unveils strong anharmonicity in non-Arrhenius self-diffusion of tungsten.
\newblock \textit{Nature Communications}, 2025. DOI: 10.1038/s41467-024-55759-w.

\bibitem[RHEA(2025)]{RHEA2025}
Utilizing a machine-learned potential to explore enhanced radiation tolerance in the MoNbTaVW high-entropy alloy.
\newblock \textit{Journal of Nuclear Materials}, 2025. DOI: 10.1016/j.jnucmat.2025.156004.

\bibitem[Tunes et~al.(2025)]{Tunes2025}
Tunes M A, Parkison D, Sun B, et~al.
\newblock High radiation resistance in the binary W-Ta system through small V additions.
\newblock \textit{Advanced Science}, 2025. DOI: 10.1002/advs.202417659.

\bibitem[Gesson et~al.(2025)]{Gesson2025}
Gesson L, Henning G, Collin J, Vanstalle M.
\newblock Enhancing nuclear cross-section predictions with deep learning: the DINo algorithm.
\newblock \textit{European Physical Journal Plus}, 2025. DOI: 10.1140/epjp/s13360-025-06562-z.

\bibitem[H-W(2025)]{HW2025}
Sticking, reflection, and abstraction behavior of hydrogen on tungsten surfaces using ML potential.
\newblock \textit{Acta Materialia}, 2025. DOI: 10.1016/j.actamat.2025.121306.

\bibitem[TMAP8(2026)]{TMAP82026}
Multiscale assessment of tritium behavior in fusion pilot plant design using surrogate models in TMAP8.
\newblock \textit{ArXiv}, 2026.

\bibitem[ITER-Gamma(2025)]{ITERGamma2025}
A machine learning case study in nuclear fusion: DT fusion power assessment with gamma-ray spectroscopy.
\newblock \textit{Energy and AI}, 2025. DOI: 10.1016/j.egyai.2025.100526.

\bibitem[Rothstein et~al.(2025)]{Rothstein2025PACMAN}
Rothstein A, Farre-Kaga H J, Butt J, et~al.
\newblock Enabling integrated AI control on DIII-D: PACMAN.
\newblock \textit{arXiv:2511.08818}, 2025.

\bibitem[Sonker et~al.(2026)]{Sonker2026}
Sonker R, Kaga H J F, Chen J, et~al.
\newblock Offline reinforcement learning for rotation profile control on DIII-D.
\newblock \textit{arXiv:2605.05857}, 2026.

\bibitem[Wu et~al.(2025)]{Wu2025}
Wu N, Li R, Yang Z, et~al.
\newblock Plasma shape control via zero-shot generative reinforcement learning.
\newblock \textit{arXiv:2510.17531}, 2025.

\bibitem[Sorokin et~al.(2026)]{Sorokin2026}
Sorokin D, Stokolesov M, Granovskiy A, et~al.
\newblock Dynamic plasma shape control with arbitrary sensor subsets.
\newblock \textit{arXiv:2605.15935}, 2026.

\bibitem[Gupta et~al.(2025)]{Gupta2025}
Gupta A, Eldon D, Bang E, et~al.
\newblock Detachment control in KSTAR with tungsten divertor.
\newblock \textit{arXiv:2505.07978}, 2025.

\bibitem[DivControlNN(2025)]{DivControlNN2025}
DivControlNN: Latent space mapping for divertor plasma detachment control.
\newblock \textit{arXiv:2502.19654}, 2025.

\bibitem[Liu \& Stacey(2025)]{Liu2025}
Liu Z, Stacey W M.
\newblock Sensitivity analysis in NeuralPlasmaODE for ITER burning plasmas.
\newblock \textit{arXiv:2507.09432}, 2025.

\bibitem[PanoMHD(2026)]{PanoMHD2026}
PanoMHD: Multimodal modelling of plasma dynamics towards tokamak control.
\newblock \textit{arXiv:2603.02672}, 2026.

\bibitem[Wan et~al.(2026)]{Wan2026WEST}
Wan C, Almuhisen F, Moreau P, et~al.
\newblock Transformer-based prediction of global plasma parameters on WEST.
\newblock \textit{arXiv:2602.19110}, 2026.

\bibitem[Mouchamps et~al.(2025)]{Mouchamps2025}
Mouchamps A, Malherbe A, Bolland A, Ernst D.
\newblock Gym-TORAX: Open-source RL environment for plasma control.
\newblock \textit{arXiv:2510.11283}, 2025.

\bibitem[Ding et~al.(2025)]{Ding2025}
Ding S, Zhang Z, Shi G, et~al.
\newblock Physics-informed neural operator for nonlinear Grad-Shafranov equation.
\newblock \textit{arXiv:2511.19114}, 2025.

\bibitem[Subbotin et~al.(2025)]{Subbotin2025}
Subbotin G F, Sorokin D I, Nurgaliev M R, et~al.
\newblock First application of deep RL for magnetic plasma control on DIII-D.
\newblock \textit{arXiv:2506.13267}, 2025.

\bibitem[Lee et~al.(2025)]{Lee2025Nature}
Lee J, et~al.
\newblock Deep learning to control plasma instabilities in tokamaks.
\newblock \textit{Nature}, 2025. DOI: 10.1038/s41586-025-08699-4.

\bibitem[Pfau et~al.(2025)]{Pfau2025Nature}
Pfau D, et~al.
\newblock Accelerating magnetic confinement fusion research with deep reinforcement learning.
\newblock \textit{Nature}, 2025. DOI: 10.1038/s41586-025-08737-1.

\bibitem[Conlin et~al.(2024)]{Conlin2024Nature}
Conlin R, et~al.
\newblock Optimizing stellarators with differentiable programming.
\newblock \textit{Nature}, 2024. DOI: 10.1038/s41586-024-07648-x.

\bibitem[Poels et~al.(2025a)]{Poels2025VAE}
Poels Y, et~al.
\newblock Plasma state monitoring and disruption characterization using multimodal VAEs.
\newblock \textit{Nuclear Fusion}, 2025. DOI: 10.1088/1741-4326/adf121.

\bibitem[Poels et~al.(2025b)]{Poels2025Class}
Poels Y, et~al.
\newblock Robust confinement state classification with uncertainty quantification.
\newblock \textit{Nuclear Fusion}, 2025. DOI: 10.1088/1741-4326/adf349.

\bibitem[Bandyopadhyay et~al.(2025)]{Bandyopadhyay2025}
Bandyopadhyay I, et~al.
\newblock MHD, disruptions and control physics: Chapter 4.
\newblock \textit{Nuclear Fusion}, 65:103001, 2025. DOI: 10.1088/1741-4326/ade7a0.

\bibitem[Carey et~al.(2025)]{Carey2025}
Carey N, et~al.
\newblock Neural operator surrogate models of plasma edge simulations.
\newblock \textit{Nuclear Fusion}, 2025. DOI: 10.1088/1741-4326/adfdfb.

\bibitem[Zheng et~al.(2025)]{Zheng2025}
Zheng G H, et~al.
\newblock EFIT-mini: an embedded, multi-task neural network equilibrium inversion.
\newblock \textit{Nuclear Fusion}, 2025. DOI: 10.1088/1741-4326/adff94.

\bibitem[Ling et~al.(2025)]{Ling2025}
Ling Y, et~al.
\newblock PaMMA-net: plasma magnetic measurement evolution.
\newblock \textit{Nuclear Fusion}, 2025. DOI: 10.1088/1741-4326/ae0655.

\bibitem[Muraca et~al.(2025)]{Muraca2025}
Muraca M, et~al.
\newblock Integrated modeling of SPARC H-mode scenarios.
\newblock \textit{Nuclear Fusion}, 2025. DOI: 10.1088/1741-4326/adf656.

\bibitem[Morosohk et~al.(2025)]{Morosohk2025}
Morosohk S, et~al.
\newblock Experimental demonstration of real-time electron temperature profile control in DIII-D.
\newblock \textit{Nuclear Fusion}, 2025. DOI: 10.1088/1741-4326/adf456.

\bibitem[Arnaud et~al.(2025)]{Arnaud2025}
Arnaud J S, et~al.
\newblock A runaway electron avalanche surrogate for partially ionized plasmas.
\newblock \textit{Nuclear Fusion}, 2025. DOI: 10.1088/1741-4326/ae00db.

\bibitem[Garcia et~al.(2026)]{Garcia2026}
Garcia J, et~al.
\newblock Overview of first JT-60SA plasma operation.
\newblock \textit{Nuclear Fusion}, 2026. DOI: 10.1088/1741-4326/ae74e1.

\bibitem[Luo et~al.(2025)]{Luo2025}
Luo Y, et~al.
\newblock Neural network for edge transport modeling on W7-X.
\newblock \textit{Nuclear Fusion}, 2025. DOI: 10.1088/1741-4326/adf75f.

\bibitem[Dudt et~al.(2024)]{Dudt2024}
Dudt D, Conlin R, Panici D, Kolemen E.
\newblock Optimization of nonlinear turbulence in stellarators.
\newblock \textit{Journal of Plasma Physics}, 2024.

\bibitem[Unalmis et~al.(2024)]{Unalmis2024}
Unalmis K E, Gaur R, Conlin R, et~al.
\newblock Differentiable bounce-averaging algorithm for stellarators.
\newblock \textit{arXiv:2412.01724}, 2024.

\bibitem[Padidar et~al.(2025)]{Padidar2025}
Padidar M, Huang T, Giuliani A, Spivak M.
\newblock Diffusion for Fusion: Designing Stellarators with Generative AI.
\newblock \textit{arXiv:2511.20445}, 2025.

\bibitem[Kaptanoglu and Gil(2026)]{Kaptanoglu2026}
Kaptanoglu A A, Gil P F.
\newblock Automated AI-driven stellarator coil optimization.
\newblock \textit{arXiv:2603.15240}, 2026.

\bibitem[Thun et~al.(2026)]{Thun2026}
Thun T, Merlo A, Conlin R, Panici D.
\newblock Improving ideal MHD equilibrium accuracy with PINNs.
\newblock \textit{Nuclear Fusion}, 2026. DOI: 10.1088/1741-4326/ae2937.

\bibitem[Angelis et~al.(2026)]{Angelis2026}
Angelis D, Sofos F, et~al.
\newblock Prediction of neutral gas pressure in W7-X.
\newblock \textit{Physics of Plasmas}, 33(1):012501, 2026.

\bibitem[Cadena et~al.(2025)]{Cadena2025}
Cadena S, Merlo A, Laude E, et~al.
\newblock ConStellaration: A dataset of QI-like stellarator boundaries.
\newblock \textit{NeurIPS Datasets and Benchmarks}, 2025.

\bibitem[Subbotin et~al.(2026)]{Subbotin2026}
Subbotin G F, Sorokin D I, et~al.
\newblock Reconstruction-free magnetic control of DIII-D plasma with DRL.
\newblock \textit{Nuclear Fusion}, 2026. DOI: 10.1088/1741-4326/ae34c6.

\bibitem[Zhang et~al.(2026)]{Zhang2026Boundary}
Zhang Q, Li T, et~al.
\newblock Deep-learning real-time plasma boundary detection on EAST.
\newblock \textit{Nuclear Fusion}, 66(3):036048, 2026.

\bibitem[Grelier et~al.(2025)]{Grelier2025}
Grelier E, Gorse V, Mitteau R, et~al.
\newblock Deep learning for WEST first wall monitoring.
\newblock \textit{IEEE Trans. Plasma Sci.}, 2025.

\bibitem[Dasbach et~al.(2026)]{Dasbach2026}
Dasbach S, Brezinsek S, et~al.
\newblock SOLPS-NN: Deep-learning surrogates for plasma exhaust.
\newblock \textit{arXiv:2604.19223}, 2026.

\bibitem[Wiesen et~al.(2024)]{Wiesen2024}
Wiesen S, Dasbach S, et~al.
\newblock Data-driven models in fusion exhaust.
\newblock \textit{Nuclear Fusion}, 64(8):086046, 2024.

\bibitem[Zhu et~al.(2025)]{Zhu2025Latent}
Zhu B, Zhao M, Xu X Q, et~al.
\newblock Latent space mapping for divertor detachment control.
\newblock \textit{Physics of Plasmas}, 32(6):062508, 2025.

\bibitem[Zhang et~al.(2025)]{Zhang2025Neutral}
Zhang J, Mao S, et~al.
\newblock Neutral source terms with deep learning for edge plasma.
\newblock \textit{Plasma Sci. Technol.}, 27(7):075106, 2025.

\bibitem[Boschi et~al.(2026)]{Boschi2026}
Boschi T, Loreti A, et~al.
\newblock TokaMind: Multi-Modal Transformer Foundation Model.
\newblock \textit{arXiv:2602.15084}, 2026.

\bibitem[Gorse et~al.(2025)]{Gorse2025}
Gorse V, Mitteau R, Marot J.
\newblock Multimodal LLM for WEST first wall monitoring.
\newblock \textit{Knowledge-Based Systems}, 2025.

\bibitem[Joglekar et~al.(2026)]{Joglekar2026}
Joglekar A S, Thomas A G R, et~al.
\newblock Differentiable programming for plasma physics.
\newblock \textit{arXiv:2603.11231}, 2026.

\bibitem[Faraji et~al.(2025)]{Faraji2025}
Faraji F, Reza M, Knoll A.
\newblock Discovery of discretized differential equations from data.
\newblock \textit{Journal of Applied Physics}, 2025.

\bibitem[Bonalumi et~al.(2024)]{Bonalumi2024}
Bonalumi D, et~al.
\newblock XAI applied to disruption prediction in tokamaks.
\newblock \textit{Frontiers in Physics}, 2024. DOI: 10.3389/fphy.2024.1359656.

\bibitem[Roy et~al.(2026)]{Roy2026Adversarial}
Roy A, et~al.
\newblock Adversarial vulnerabilities in neural operator digital twins.
\newblock \textit{arXiv}, 2026.

\bibitem[Agnello et~al.(2026)]{Agnello2026}
Agnello A, et~al.
\newblock Challenges and opportunities for AI in fusion energy.
\newblock \textit{arXiv:2603.25777}, 2026.

\bibitem[Gahle and Barbarino(2026)]{Gahle2026}
Gahle D S, Barbarino M.
\newblock IAEA Fusion Data Lake Project.
\newblock \textit{arXiv:2604.01797}, 2026.

\bibitem[Pankin et~al.(2026)]{Pankin2026}
Pankin A Y, et~al.
\newblock NIMROD-to-IMAS workflow for extended-MHD data.
\newblock \textit{arXiv:2605.23121}, 2026.

\bibitem[Citrin et~al.(2024)]{Citrin2024TORAX}
Citrin J, et~al.
\newblock TORAX: A Fast and Differentiable Tokamak Transport Simulator.
\newblock \textit{arXiv:2406.06718}, 2024.

\end{thebibliography}

\end{CJK}
\end{document}
